\scanpage{274}
\mspage{438}

and the equation
\[
\xi=\eta\varepsilon_1(\mu-1)
\]
\deleted{\unsure{is equivalent to}} implies
\[
(52)\qquad B=S,\quad D=U,
\]
and in fact it \unsure{is equivalent} to one of these two equations and to
\[
(53)\qquad \mu D=1\qquad (\unsure{\Longleftrightarrow|\Longrightarrow}\ \mu U=1)
\]
when we already assume that $\xi,\eta$ \unsure{satisfy} the equation of the \deleted{traces and the} determinants \inserted{and the traces}
\[
(54)\quad
\left\{\begin{array}{ll}
AD-BC=1 & RU-ST=1\\
\unsure{C|A}+D-B=1 & T-S=1
\end{array}\right.
\]

\unsure{The} \deleted{\illegible{1}} scheme of solutions of the simultaneous equations $(33)$ \unsure{to|:} $(35)$ is then that of the $(D,B)$ with $D$ invertible, or again of the $(S,U)$ with $S$ invertible. We can identify it with the scheme of matrices
\[
(55)\qquad
\begin{pmatrix}\mu&M\\0&1\end{pmatrix}
=\begin{pmatrix}1/D&B/D\\0&1\end{pmatrix}
\]
(cf. page~410 \illegible{1}), which \unsure{normalize} the $1$-\unsure{parameter} subgroup
\[
\deleted{\unsure{\rho}\,\unsure{\varepsilon}}\quad
\varepsilon_0(\unsure{\lambda})=\deleted{\illegible{1}}
\begin{pmatrix}1&-\lambda\\0&1\end{pmatrix}
=\rho^{-1}\varepsilon_1(\lambda)\rho.
\]


\scanpage{275}

Let us return to the general case, where we nevertheless assume that $\varepsilon_1$ given by $(40\ \mathrm{bis})$, \deleted{\unsure{then}} let us put
\[
(56)\qquad \varepsilon_0=\rho^{-1}\varepsilon_1\rho
=-\lambda_0\sigma^{-1}\rho=1+\unsure{N_0},\quad\text{with}
\]
\[
(57)\qquad N_0=\rho^{-1}N_1\rho=-(\lambda_0\sigma^{-1}\rho+1),
\]
so that $N_0^2=0$, and
\[
(58)\qquad \varepsilon_0\quad\lambda\longmapsto
\varepsilon_0(\lambda)\overset{\unsure{\mathrm{dfn}}}{=}1+\lambda N_0
\]
is a $1$-\unsure{parameter} subgroup of $SG(\mathcal M)$, say $L_0$, whose \unsure{normalizer} in $G(\mathcal M)$ will be denoted by $B_0$. We have two \unsure{homomorphisms}

$(59)$
\diagram{\begin{tikzcd}[ampersand replacement=\&]
B_0 \arrow[r, shift left=1ex, "\chi(\underline u)=\operatorname{dét}(\underline u)"]
\arrow[r, shift right=1ex, "\chi'(\underline u)\ \text{given}"'] \& \unsure{\mathbf G_m|G_m}
\end{tikzcd}}
by
\[
\underline u(\varepsilon_0(\lambda))=\varepsilon_0(\chi'(\underline u)\lambda)
\]
or again (in linear terms in $\mathcal M$)
\[
(60)\qquad \underline uN_0\underline u^{-1}=\chi'(\underline u)N_0
\]
\margin{I underline $\underline u$, so as not to confuse it with a coefficient of $\unsure{\eta|\beta}=\begin{pmatrix}r&s\\t&u\end{pmatrix}$ !}
And we \unsure{consider} the subgroup
\[
(61)\quad\left\{\begin{array}{l}
\widetilde B_0^{\unsure{*}}\subset B_0\quad\text{defined by }\chi(\underline u)=\chi'(\underline u)\\
\text{i.e. by }\deleted{\unsure{\det}}\\
\underline uN_0\underline u^{-1}=\det(\underline u)N_0
\end{array}\right.
\]

so we have
\[
L_0\subset\widetilde B_0^*.
\]
and we even have an exact sequence of group schemes
\diagram{\begin{tikzcd}[ampersand replacement=\&]
1\arrow[r]\&L_0\arrow[r]\&\widetilde B_0^*\arrow[r,"\det"]\&\unsure{\mathbf G_m|G_m}\arrow[r]\&1.
\end{tikzcd}}
\margin{we would need to \unsure{check} \illegible{4}\ \deleted{\illegible{1}}\quad\unsure{This is} \uline{false}, \unsure{cf. below}.}


\scanpage{276}
\mspage{439}

Let $\underline u$ be a section of $\widetilde B_0^*$, I feel like proving that if we put
\[
(63)\qquad \xi=[\underline u,\rho],\quad\eta=[\underline u,\sigma],
\]
we will automatically have the relation
\[
(64)\qquad\xi=\eta\varepsilon_1(\mu-1)\qquad\text{where }\mu=\det\underline u.
\]
To \unsure{reassure} ourselves, let us first take $\underline u$ in $L_0$, i.e.
\[
\underline u=\varepsilon_0(\lambda)
\]
whence
\[
\begin{aligned}
\relax
[\underline u,\rho]&=\varepsilon_0(\lambda)\rho\varepsilon_0(\lambda)^{-1}\rho^{-1}
=\varepsilon_0(\lambda)\varepsilon_1(-\lambda)\\
[\underline u,\sigma]&=\varepsilon_0(\lambda)\sigma\varepsilon_0(-\lambda)\sigma^{-1}
\end{aligned}
\]
Now
\[
\sigma\underbrace{\varepsilon_0(-\lambda)}_{\inserted{1+N_0\lambda}}\sigma^{-1}
=1+(\sigma\uline{N_0}\sigma^{-1})\lambda\ \deleted{\unsure{=}}
\]
\unsure{now} \deleted{\illegible{1}} by virtue of $(57)$ we have
\[
\sigma N_0\sigma^{-1}=-(\lambda_0\rho\sigma^{-1}+1)\overset{(49)}{=}N_1
\]
so we have
\[
(65)\quad\left\{\begin{array}{ll}
\sigma\varepsilon_0(\lambda)\sigma^{-1}=\varepsilon_1(\lambda)&\text{and similarly}\\
\rho\varepsilon_0(\lambda)\rho^{-1}=\varepsilon_1(\lambda)
\end{array}\right.
\]
so we do indeed find
\[
[\underline u,\rho]=[\underline u,\sigma]=\varepsilon_0(\lambda)\varepsilon_1(-\lambda)
\]
i.e. $(64)$, since here $\unsure{\mu}=\det u=\unsure{1}$, so $\mu-1=0$.


\scanpage{277}

We have seen that $\rho$ and $\sigma$ both \unsure{transform} $N_0$ into $N_1$, hence the additive $1$-\unsure{parameter} subgroup
\[
L_0=\left\{\underset{\inserted{=1+\lambda N_0}}{\unsure{\varepsilon_0}(\lambda)}\ \middle|\ \lambda\text{ in }\unsure{\mathbf G_a|G_a}\right\}
\]
into the analogous subgroup
\[
L_1=\left\{\underset{\inserted{=1+\lambda N_1}}{\varepsilon_1(\lambda)}\ \middle|\ \lambda\text{ in }\unsure{\mathbf G_a|G_a}\right\},
\]
and hence $\widetilde B_0$ into $\widetilde B_1$ (defined in terms of $L_1$, or of $N_1$, as $\widetilde B_0^*$ in terms of $L_0$). \deleted{Thus} The formula $(\unsure{64|69})$ can also be written
\margin{\deleted{$\rho^{-1}\sigma$}}
\[
\rho\underline u^{-1}\rho^{-1}=\unsure{\sigma}\underline u^{-1}\sigma^{-1}\varepsilon_1(\mu-1)
\]
or also, replacing $\underline u$ by $\underline u^{-1}$
\[
(\unsure{*|**})\qquad
\rho\underline u\rho^{-1}=\sigma\underline u\sigma^{-1}\varepsilon_1\left(\frac1\mu-1\right).
\]
Now we \deleted{\illegible{1}} have
\[
\rho\underline u\rho^{-1}=\deleted{\unsure{\lambda}}\operatorname{int}(\lambda)(\sigma\underline u\sigma^{-1})
\]
where $\uline{\lambda=\rho\sigma^{-1}}$, and the formula $\unsure{(*)|(**)}$ can also be written
\[
(**):\qquad
\operatorname{int}(\lambda)(v)=v\cdot\varepsilon_1\left(\frac1\mu-1\right)
\qquad\begin{gathered}
\text{for }v\ \uline{\text{in }\unsure{B'_1|\widetilde B_1}}\\
\text{and }\mu=\det v;
\end{gathered}
\]
and by virtue of $(40\ \mathrm{bis})$ $\operatorname{int}\lambda=\operatorname{int}(\varepsilon_1)$, so the formula reduces to
\[
\operatorname{int}(\varepsilon_1)(v)=v\cdot\varepsilon_1\left(\frac1\mu-1\right)
\]


\scanpage{278}
\mspage{440}

or again
\[
\deleted{(1+N_1)\unsure{v_0|v}(1-N_1)=v\left(1+N_1\left(\frac1\mu-1\right)\right)}
\]
\deleted{i.e.}
\[
\deleted{v+N_1v-vN_1\unsure{+|-}N_1vN_1=v+vN_1\left(\frac1\mu-1\right)}
\]
\[
\varepsilon_1v\varepsilon_1^{\deleted{\unsure{-1}}}(-1)=v\varepsilon_1\left(\frac1\mu-1\right)\quad\text{i.e.}
\]
\[
\underbrace{\varepsilon_1}_{\inserted{1+N_1}}v
=v\underbrace{\varepsilon_1\left(\frac1\mu\right)}_{\inserted{1+\frac1\mu N_1}}
\qquad\text{i.e. }v+N_1v=v+\frac1\mu vN_1
\]
that is, finally
\[
vN_1v^{-1}=\mu N_1\quad(=\det(v).N_1)
\]
which is nothing other than the definition of $\widetilde B_1^{\unsure{*}}$ !

We thus find a \deleted{\unsure{mono}}morphism
\[
(66)\quad\left\{\begin{array}{rcl}
\widetilde B_0&\longrightarrow&\mathcal X_{\rho,\sigma}\\
\underline u&\longmapsto&([\underline u,\rho],[\underline u,\sigma],\det\underline u)
\end{array}\right.
\]
from $\widetilde B_0^{\unsure{*}}$ into the set of solutions of the system $(33)$ to $(35)$, is it an isomorphism?

\unsure{Note} that the equality $[\underline u,\rho]=[\underline u',\rho]$ means that $u^{-1}\unsure{u'}\in\Gamma T_\rho$ so \deleted{\illegible{3}} saying that $(66)$ is a monomorphism simply means that
\[
\deleted{(\unsure{67})}\qquad T_\rho\cap\widetilde B_0=\{1\}
\]
in fact the analysis carried out must \unsure{show}

\scanpage{279}

that this intersection is also equal to $T_\sigma\cap\widetilde B_0^{\unsure{*|\prime}}$, hence \unsure{contained} in $T_\sigma\cap T_\rho$. I am going to assume
\[
\tag{68}
T_\sigma\cap T_\rho=\mathbf G_m\;(=\unsure{\mathrm{Cent}}\,G(M))
\]
\inserted{(\unsure{is} $\unsure{\supset|\supseteq}$, it \illegible{2} $\unsure{\subset|=}$)}
--- it \unsure{will suffice|suffices} for this to assume that in each fiber, $1,\sigma,\rho$ are linearly independent --- this is of course the case in the case we are interested in. \uline{But} Thus
\[
T_\sigma\cap T_\rho\cap\widetilde B_0^{\unsure{*|\prime}}\subset\mathbf G_m\cap\widetilde B_0^*,
\]
but
\[
\mathbf G_m\cap\widetilde B_0^*=\deleted{\{1\}}\mu_2
\]
For if $\lambda$ \deleted{\illegible{1}} \unsure{is} a \unsure{local} section of $\mathbf G_m$, then
\[
\lambda N_0\lambda^{-1}=N_0=1.\unsure{N_0|N_3}
\]
but $\det\lambda=\lambda^2$, and the condition that $\lambda$ \deleted{\unsure{be}} (considered as a section of $B_0$) be in $B_0^{\unsure{*|\prime}}$, is indeed \uline{$\lambda^2=1$}. \deleted{\unsure{$B_0$}}

We find \deleted{\illegible{1}} in this way
\[
\tag{69}
T_\sigma\cap\widetilde B_0^{\unsure{*|\prime}}=T_\rho\cap\widetilde B_0^{\unsure{*|\prime}}=\mu_2
\]
and \unsure{the}\deleted{\illegible{1}} morphism (66) is not a monomorphism, but \uline{factors}

\scanpage{280}
\mspage{441}

as a monomorphism
\[
\tag{70}
\widetilde B_0/\mu_2\hookrightarrow\mathcal X_{\rho,\sigma}.
\]
To tell the truth, we see that $\widetilde B_0$ is not the ``right'' group to introduce here --- in the special case $\varepsilon_0=\begin{pmatrix}1&-1\\0&1\end{pmatrix}$, \deleted{\unsure{$B_0$}} we have
\[
B_0=\left\{\begin{pmatrix}p&\unsure{\ell|l}\\0&q\end{pmatrix}\mid pq\ \mathrm{inv.}\right\}
\]
and the relation \deleted{\unsure{for}} $\chi'(\underline u)=\underbrace{\chi(\underline u)}_{\inserted{\det\underline u=pq}}$ can be written, by an immediate computation
\[
\tag{71}
p/q=pq\qquad\mathrm{i.e.}\quad q^2=1
\]
whereas the ``right'' group we are interested in is defined by
\[
B'_0=\left\{\begin{pmatrix}p&\unsure{\ell|e}\\0&1\end{pmatrix}\mid p\ \mathrm{inv}\right\},
\]
\unsure{thus} with $q=1$. To give the description in the general case, we note that the determinant function on $B_0^*$ decomposes intrinsically into two factors
\[
\tag{72}
\det\underline u=p(\underline u)q(\underline u),
\]
which we can distinguish from each other \unsure{precisely} by the condition that
\[
\tag{73}
\underline uN_0\underline u^{-1}=\frac{p(\underline u)}{q(\underline u)}N_0,
\]
and the condition which is going to describe the ``right'' group for us is not (71), but rather
\[
\tag{74}
q=q(\underline u)=1,
\]

\scanpage{281}

which thus defines a group that is \uline{smooth} (even in characteristic 2 !) and this time fits into a canonical exact sequence
\[
\tag{74}
1\longrightarrow\unsure{L_0^*|L_0}\longrightarrow B'_0\xrightarrow{\det}\mathbf G_m\longrightarrow1
\]
We now have
\[
\tag{75}
B'_0\cap\mathbf G_m=\{1\},
\]
which implies that the morphism induced by (66)
\[
\tag{76}
B'_0\hookrightarrow\mathcal X_{\rho,\sigma}
\]
is a monomorphism.

To show that it is an isomorphism, it now suffices, by virtue \unsure{of general principles}, to show that it is \uline{surjective} on the points with values in an algebraically closed field (\unsure{if that} \unsure{can|for} simplify life for us). I drop the details \unsure{of the|of} verifications.

\deleted{\unsure{Remark}}

\uline{\unsure{Remark|Remark}} We can consider, by condition (40 bis), that we start from a unipotent section \inserted{$\varepsilon_1=1+N_1\quad(N_1^2=0)$} \deleted{\unsure{$\varepsilon_1$}} of $SG(M)$, and from a section $\rho$, and that $\sigma$ is deduced from it by \deleted{\unsure{a}}\inserted{the}

\scanpage{282}
\mspage{442}

formula
\[
\tag{77}
\sigma=-\varepsilon_1^{-1}\rho\;\left(=-(1-N_1)\rho\quad\mathrm{i.e.}\quad\sigma=-\rho+N_1\rho,\quad\mathrm{i.e.}\quad\rho+\sigma=N_1\rho\right)
\]

\margin{I take the constant which \unsure{occurs} \unsure{in|in} the case that \unsure{interests} us most, where $-\sigma^{-1}=\sigma$, so that (40 bis) becomes $\varepsilon_1=\rho\sigma$.}

(I drop the constant \inserted{\unsure{pure|for}} $\lambda_0$, for \deleted{\illegible{1}} the situation is not changed when we replace $\rho,\sigma$ by scalar multiples, for $F_\rho$ and $F_\sigma$ is not changed \ldots) Condition (42) then becomes
\[
c=\operatorname{Tr}\rho+\operatorname{Tr}\sigma=\operatorname{Tr}(1-\underbrace{\varepsilon_1^{-1}}_{\inserted{1-N_1}})\rho=\operatorname{Tr}N_1\rho\in\Gamma\mathcal O_S^*
\]

\margin{\uline{N.B.} we have
\[
\tag{78 bis}
\left\{\begin{aligned}
\rho+\sigma&=N_1\rho=\rho N_0\\
&=-N_1\sigma=-\sigma N_0
\end{aligned}\right.
\]
\unsure{let} $\unsure{\gamma_0|\varepsilon_0}$ \illegible{4}}

(\uline{invertible} section of $\mathcal O_S$). We assume that $1,\rho,\sigma$ are linearly independent in each fiber, it amounts to the same to say that $1,\rho$, and $N_1\rho$ are, or again that $\rho^{-1},1$ and $N_1$ are, or that $1,\sigma,N_1\sigma$ are, or $\sigma^{-1},1,N_1$. In summary

\margin{\unsure{N.B.} (After possibly multiplying \unsure{the|the} $\rho$ by an invertible scalar) \unsure{we} \uline{$\operatorname{Tr}\unsure{\gamma_0|\delta_0}=1$} ($\unsure{\gamma_0|\delta_0}=N_1\rho$) ! (which normalizes \unsure{entirely}, with the relation $\varepsilon_1=-\rho\sigma^{-1}$, the \unsure{arbitrary|intrinsic} scalar factors of $\rho,\sigma$)}

\uline{Proposition}. \unsure{Let} $\unsure{M|\mathcal M}$ be an Azumaya algebra over the base scheme $S$, $N_1$ a section of square zero \deleted{\unsure{which}} i.e. $\varepsilon_1=1+N_1$ a unipotent section, $\rho$, \inserted{$\sigma$} \deleted{\unsure{a}} \inserted{some} sections related by
\[
\tag{\unsure{40|78}}
\varepsilon_1=-\rho\sigma^{-1}\quad\mathrm{i.e.}\quad\sigma=-\varepsilon_1^{-1}\rho\quad\mathrm{i.e.}\quad\rho+\sigma=N_1\rho\quad\mathrm{i.e.}\quad\unsure{\sigma+\rho|\rho+\sigma}=-N_1\sigma
\]
\deleted{\illegible{1}} (which thus determines each of the three sections \unsure{\deleted{$\varepsilon_1$}|$\varepsilon_1$} $\rho,\sigma$ and $\varepsilon_1$ (or $N_1$) in terms of the other two, and we can choose arbitrarily $\rho$ and \deleted{\unsure{$N_1$}} $\varepsilon_1$ i.e. $\rho$ and $N_1$ of square zero) --- we